\section{Practical guidance and limits}
\label{sec:discussion}

\paragraph*{Guidance.}
(1)~On the CLD flow, an isotropic schedule needs $r\geq\kappa$ for the standard clipped reverse SDE.
Below $\kappa$, use the corrected sampler \Sb{}: it needs no retraining, and it removes the threshold effect with the exact and the learned score alike.
(2)~Below $\kappa$, the damage is set by $\smin$ relative to the data scale and by the mass $M$, not by $\smax$.
With $\smin$ well below the data scale and the CLD mass, it stays at a few per cent.
(3)~With a learned score, prefer the corrected SDE to the probability-flow ODE: the ODE carries score errors instead of forgetting them.

\paragraph*{Limits.}
The standard CLD, whose noise is given by its SDE, does not meet this threshold.
In typical settings ($\smin$ below the data scale, setting D) the threshold effect is a few per cent, smaller than the error of a small learned score.
The large effects (11--19\% in mixtures) require $\smin$ comparable to the data scale or a mass far from the default.
All data are four-dimensional Gaussian mixtures; images were not attempted, because the map predicts the same conclusion there.
Three observations remain open: the velocity effect in setting A, the mechanism by which neighbouring modes reduce the shrinkage, and the small ODE-to-SDE ratio for the final networks of setting C.

\begin{conjecture}
For phase-space samplers whose drift has a skew part, such as the CLD flow, the fragility of the probability-flow ODE with a learned score (Observation~\ref{obs:ode}) is a consequence of Eq.~\eqref{eq:M}, independent of the threshold, wherever the covariance is dominated by the blur.
\end{conjecture}
